\section{Graded Wheel Setup and Category Discipline}
\label{sec:setup}

We first define the arithmetic source named in the title, then isolate the
smaller set of hypotheses used by each obstruction.

\subsection{The frozen wheel-sieve source}

Set $Q_0=1$, $q_1=2$, and $Q_1=2$.  For $k\geq1$, define
\begin{equation}
  q_{k+1}=\min\{n>q_k:\gcd(n,Q_k)=1\},
  \qquad Q_{k+1}=Q_kq_{k+1}.
  \label{eq:wheel-recurrence}
\end{equation}
This recurrence contains no prime table.

\begin{lemma}[Prime enumeration]
\label{lem:prime-enumeration}
If $p_k$ denotes the $k$th rational prime, then
$q_k=p_k$ and $Q_k=\prod_{j=1}^k p_j$ for every $k\geq1$.
\end{lemma}

\begin{proof}
The assertion holds at $k=1$.  Assume it through level $k$.  The next prime
$p_{k+1}$ is coprime to $Q_k$, so the minimum in
\eqref{eq:wheel-recurrence} exists and $q_{k+1}\leq p_{k+1}$.  If its
minimizer were composite, let $r<q_{k+1}$ be a prime divisor.  If
$r\leq p_k$, then $r\mid Q_k$, contradicting
$\gcd(q_{k+1},Q_k)=1$.  If $r>p_k=q_k$, then
$\gcd(r,Q_k)=1$, so $r$ is a smaller admissible integer than the minimizer,
again a contradiction.  Hence $q_{k+1}$ is prime.  Minimality then makes it
the first prime after $p_k$, namely $p_{k+1}$, and the formula for $Q_{k+1}$
follows by induction.
\end{proof}

At level $k\geq0$, let
\begin{equation}
  R_k=\{0\leq r<Q_k:\gcd(r,Q_k)=1\}.
\end{equation}
For $r\in R_k$, put an edge to each residue
$r+jQ_k\in R_{k+1}$ with $0\leq j<q_{k+1}$ except the unique choice
divisible by $q_{k+1}$.  Let $X_k$ be the one-sided tail paths whose first
edge begins at level $k$, and let $\sigma$ delete the first edge.  Thus
\begin{equation}
  X=\bigsqcup_{k\geq 0} X_k,
  \qquad
  \sigma(X_k)\subseteq X_{k+1}.
  \label{eq:graded-space}
\end{equation}
The exact transition clock is generated by the same recurrence:
\begin{equation}
  \tau_k=\log\frac{Q_{k+1}}{Q_k}=\log q_{k+1}.
  \label{eq:wheel-clock}
\end{equation}

\subsection{Abstract hypotheses and project convention}

\begin{definition}[Strictly graded shift]
Let $X=\bigsqcup_{k\geq0}X_k$ be a disjoint union and let
$\sigma:X\to X$ satisfy
\begin{equation}
  \sigma(X_k)\subseteq X_{k+1}
  \qquad(k\geq0).
  \label{eq:strict-grading}
\end{equation}
We call $(X,\sigma)$ a strictly graded shift and write
$\level(x)=k$ for $x\in X_k$.
\end{definition}

The wheel system above is strictly graded.  The quotient result below uses
only that $q_{k+1}$ is distinct across levels; the decoder result uses only
that the required output set is infinite.

\begin{definition}[Level-blind stationarization convention]
\label{def:stationarization}
For this project, a \emph{stationarization proposal} specifies one target
phase space $Y$, one self-map $S:Y\to Y$, one declared relation to the source,
and total arithmetic and clock decoders.  The proposal is \emph{level-blind}
when the target rule and decoders are fixed independently of the source level
$k$ and every finite cutoff $K$, and their inputs contain neither an external
level nor a stored table of $(q_k)$ or $(Q_k)$.  This is an admissibility
convention, not an assertion that such a target exists.
\end{definition}

\begin{definition}[Strict extension]
A wheel-preserving strict extension is a dynamical system $(Y,S)$ together
with a total map $\pi:Y\to X$ satisfying
\begin{equation}
  \pi\circ S=\sigma\circ\pi.
  \label{eq:semiconjugacy}
\end{equation}
Only this equivariance relation is used below; surjectivity of $\pi$ is not
assumed.  This convention avoids relying on competing uses of the word
\emph{semiconjugacy}.
\end{definition}

For graph quotients, let $G=(V,E)$ be a directed graph and let $\sim$ be an
equivalence relation on $V$.

\begin{definition}[Strong forward bisimulation]
The relation $\sim$ is a strong forward bisimulation when, for all
$v\sim w$,
\begin{enumerate}[label=(\roman*)]
  \item if $vEv'$, then some $w'$ satisfies $wEw'$ and $v'\sim w'$; and
  \item if $wEw'$, then some $v'$ satisfies $vEv'$ and $v'\sim w'$.
\end{enumerate}
The quotient graph $G/{\sim}$ has an edge $[v]\to[v']$ exactly when some
representatives $u\in[v]$ and $u'\in[v']$ satisfy $uEu'$.  The obstruction
uses this definition only through successor matching.
\end{definition}

\begin{definition}[Forward well-founded graph]
\label{def:forward-well-founded}
A directed graph is forward well-founded when it has no infinite directed
path $v_0Ev_1Ev_2E\cdots$.
\end{definition}

Finally, a finite-local observation consists of a finite alphabet $A$, a
fixed finite coordinate window $W$, and one decoder $d:A^W\to Z$ into a
declared codomain $Z$.  The words \emph{fixed} and \emph{finite} exclude a
cutoff-dependent lookup table or a window that grows with the level.

\begin{table}[H]
\centering
\small
\caption{Each obstruction is tied to an explicit hypothesis class.  The
rightmost column is part of the result, not a proposed construction.}
\label{tab:scope}
\begin{tabularx}{\textwidth}{@{}>{\raggedright\arraybackslash}p{0.22\textwidth}
  >{\raggedright\arraybackslash}p{0.29\textwidth}
  >{\raggedright\arraybackslash}X@{}}
\toprule
Class & Conclusion & Not covered \\
\midrule
Strict extension satisfying \eqref{eq:semiconjugacy} &
Target has no periodic points; independently, the source's full-backward
inverse limit is empty &
New factor, quotient, recoding, or independent grammar \\
Forward-well-founded graph plus strong forward bisimulation &
Quotient acyclic &
Infinite graph with an infinite forward path; weaker graph image \\
Representative-exact, level-injective state-class decoder &
Quotient inherits strict grading &
Edge or path decoder that distinguishes representatives \\
Finite alphabet plus fixed finite window &
Infinite exact output set undecodable &
Countable alphabet, infinite memory, or another unbounded input \\
\bottomrule
\end{tabularx}
\end{table}

The discipline in \cref{tab:scope} prevents two errors.  First, a theorem
cannot be enlarged by silently dropping one of its hypotheses.  Second, an
object outside a no-go theorem does not receive positive evidence merely by
escaping that theorem.
